
% ==============================================================================
% SECTION 2: CALCULS D'AUTONOMIE
% ==============================================================================
\section{Calculs d'Autonomie}

\begin{frame}{Introduction}
    La planification de la plongée nécessite des calculs d'autonomie précis 
    pour garantir la sécurité de la palanquée, surtout en absence de Directeur de Plongée (DP).
    \begin{block}{Objectifs}
        \begin{itemize}
            \item Comprendre la relation entre pression, volume et consommation d'air.
            \item Savoir calculer l'autonomie en fonction de la profondeur et de la consommation.
            \item Appliquer la règle des tiers pour une gestion sécurisée du gaz.
        \end{itemize}
    \end{block}
\end{frame}

% --- D. LES BOUTEILLES/ÉPROUVETTES (Volumes) ---
% Macro pour dessiner une bouteille
% #1: x, #2: y, #3: Hauteur remplissage (0 à 1), #4: Label Volume, #5: Label Math
\newcommand{\drawbottle}[5]{
    % #1: x, #2: y
    % #3: Volume d'AIR (1.0 = plein, 0.5 = moitié, 0.25 = quart)
    % #4: Texte Volume, #5: Texte Math
    
    \begin{scope}[shift={(#1,#2)}]
        % --- 1. LE REMPLISSAGE (L'EAU D'ABORD) ---
        % On remplit TOUT le volume possible en EAU (Cyan)
        % Le goulot (en bas) + le corps (en haut)
        % Comme ça, le bas ne sera jamais blanc.
        
        % Le Goulot (Trapèze du bas)
        \fill[cyan!40] (-0.1, -1) -- (0.1, -1) -- (0.3, -0.8) -- (-0.3, -0.8) -- cycle;
        
        % Le Corps principal (Rectangle) - On le remplit totalement d'eau par défaut
        \fill[cyan!40] (-0.3, -0.8) rectangle (0.3, 1);

        % --- 2. L'AIR (LE BLANC PAR DESSUS) ---
        % On calcule la hauteur d'AIR.
        % Hauteur totale du corps = 1.8 (de -0.8 à 1.0)
        % Hauteur Air = 1.8 * #3
        % Le rectangle d'air part du HAUT (1.0) et descend.
        \fill[white] (-0.28, 0.98) rectangle (0.28, {1.0 - (1.8 * #3)});

        % --- 3. LE CONTOUR (NOIR) ---
        % On redessine les traits noirs par dessus les couleurs
        \draw[thick, black] (-0.3, 1) -- (-0.3, -0.8) -- (-0.1, -1) -- (0.1, -1) -- (0.3, -0.8) -- (0.3, 1) -- cycle;
        % Petite courbe pour fermer le haut joliment (optionnel)
        %\draw[thick, black] (-0.3, 1) arc(180:0:0.3cm and 0.05cm); 

        % --- 4. TEXTES ---
        % Texte volume à droite
        \node[right, text=blue!80!black, font=\bfseries\small] at (0.4, 0) {#4};
        
        % Texte Math (= 12) loin à droite
        \node[right, text=black, font=\bfseries\large] at (4.5, 0) {= 12};
    \end{scope}
}

\begin{frame}{Rappels Physiques}
    \begin{columns}
        \begin{column}{0.6\textwidth}
            \centering
            \begin{tikzpicture}[
                scale=0.65, transform shape, % Ajuste la taille pour que ça rentre
                font=\sffamily\bfseries,
                % Définition des couleurs spécifiques à l'image
                colEau/.style={fill=cyan!30},
                colTerre/.style={fill=brown!60!black},
                colTexteRouge/.style={text=red, fill=white, inner sep=2pt, font=\bfseries},
                colBoite/.style={fill=white, text=blue!80!black, draw=blue!80!black, thin, inner sep=3pt},
                arrowStyle/.style={->, >=Latex, red, line width=4pt}
            ]

            % --- A. DÉCORS (Fond) ---
            % 1. L'EAU : On la fait commencer à x=0 pour qu'elle passe SOUS la falaise.
            % Plus de trou possible !
            \fill[cyan!20] (0, 0.5) rectangle (14, -7.5);
            
            % Surface de l'eau (Vagues)
            \draw[cyan!50, decorate, decoration={snake, amplitude=1.5pt, segment length=5pt}] (0, 0.5) -- (14, 0.5);

            % La Falaise (à gauche)
            \fill[brown!70!black] (-0.5,1) -- (2,1) -- (2.2, -0.5) -- (1.8, -2) -- (2, -4) -- (1.5, -7.5) -- (-0.5, -7.5) -- cycle;
            \node[text=blue!80!black, anchor=south west] at (-0.5, 0.5) {Surface};

            % --- B. LA FLÈCHE DE PRESSION (Axe central) ---
            % On trace une grande flèche rouge vers le bas
            \draw[arrowStyle] (4, 0.5) -- (4, -7);
            \node[text=red, align=center, font=\bfseries\small] at (4, 1) {P. Absolue\\1 bar};

            % --- C. LES ÉTAGES (Profondeur & Pression) ---
            
            % Niveau 10m
            \node[colBoite] at (1, -2) {10m};
            \node[colTexteRouge] at (4, -2) {2 bars};

            % Niveau 20m
            \node[colBoite] at (1, -4.5) {20m};
            \node[colTexteRouge] at (4, -4.5) {3 bars};

            % Niveau 30m
            \node[colBoite] at (1, -7) {30m};
            \node[colTexteRouge] at (4, -7) {4 bars};

            % Dessin des 4 bouteilles
            % Surface (0m) : 1 bar -> Vol 12 (100%)
            \drawbottle{7}{1}{1}{12 litres}{}
            
            % 10m : 2 bars -> Vol 6 (50%)
            \drawbottle{7}{-1.5}{0.5}{6 litres}{}
            
            % 20m : 3 bars -> Vol 4 (33%)
            \drawbottle{7}{-4}{0.33}{4 litres}{}
            
            % 30m : 4 bars -> Vol 3 (25%)
            \drawbottle{7}{-6.5}{0.25}{3 litres}{}

            \end{tikzpicture}
        \end{column}

        \begin{column}{0.35\textwidth}
            \begin{block}{Loi de Mariotte}
                Les volumes gazeux sont inversement proportionnels à la pression.
                $$ P \times V = \text{Constante} $$
                $$ P_1 \times V_1 = P_2 \times V_2 $$
            \end{block}
        \end{column}
    \end{columns}
\end{frame}

\begin{frame}{La Physique en application à l'autonomie}
    \begin{block}{Rappel}
        L'air délivré par le détendeur est ramené à la pression ambiante (ou absolue) 
        pour pouvoir être respiré sous l'eau.  
    \end{block}
    \begin{exampleblock}{Exemple de consommation (20L/min surface)}
        \begin{itemize}
            \item Surface (1 bar) : 20 L/min
            \item 10m (2 bars) : $20 \times 2 = 40$ L/min
            \item 20m (3 bars) : $20 \times 3 = 60$ L/min
        \end{itemize}
    \end{exampleblock}
    \begin{alertblock}{Attention}
        La consommation d'air augmente avec la profondeur. 
        Il est donc crucial de planifier ses plongées en tenant compte de l'autonomie.
    \end{alertblock}
\end{frame}

% --- FRAME 1 : EXERCICE (À afficher avant pour les élèves) ---
\begin{frame}{Exercice : Calcul d'Autonomie}
    \textbf{Données :}
    \begin{itemize}
        \item Plongeur avec un Bloc de \textbf{15 Litres} gonflé à \textbf{200 bars}.
        \item Consommation en surface estimée à \textbf{20 L/min}.
        \item On néglige la réserve pour cet exercice (Volume dispo $\approx$ \textbf{?} L).
    \end{itemize}
    
    \vspace{0.5cm}
    
    \textbf{À vous de jouer : complétez le tableau suivant.}
    
    \begin{table}[]
        \centering
        \renewcommand{\arraystretch}{1.5} % Agrandit les cases pour la lisibilité
        \begin{tabular}{cccc}
        \toprule
        \textbf{Prof.} & \textbf{Pression (Pabs)} & \textbf{Conso à Prof.} & \textbf{Autonomie} \\
        \midrule
        0 m    & 1 bar       & 20 L/min         & \textbf{?} \\ 
        \hline
        20 m   & \textbf{?}  & \textbf{?}       & \textbf{?} \\ 
        \hline
        40 m   & \textbf{?}  & \textbf{?}       & \textbf{?} \\ 
        \bottomrule
        \end{tabular}
    \end{table}
\end{frame}

% --- FRAME 2 : SOLUTION (Votre frame originale améliorée) ---
\begin{frame}{Solution : Calcul d'Autonomie}
    \begin{columns}
        \begin{column}{0.4\textwidth}
            $$ \text{Autonomie} = \frac{\text{Air disponible}}{\text{Consommation à prof.}} $$
        \end{column}
        \begin{column}{0.6\textwidth}
            \begin{exampleblock}{Rappel Données}
                Bloc 15L à 200b ($\approx$ 3000L). Conso surface : 20 L/min.
            \end{exampleblock}
        \end{column}
    \end{columns}

    \begin{table}[]
        \centering
        \renewcommand{\arraystretch}{1.3}
        \begin{tabular}{cccc}
        \toprule
        Prof. & Pression & Conso ($P \times 20$) & Autonomie \\
        \midrule
        0m    & 1b       & 20 L/min     & $3000 / 20 = \mathbf{150 \text{ min}}$ \\
        20m   & 3b       & 60 L/min     & $3000 / 60 = \mathbf{50 \text{ min}}$ \\
        40m   & 5b       & 100 L/min    & $3000 / 100 = \mathbf{30 \text{ min}}$ \\
        \bottomrule
        \end{tabular}
    \end{table}
    
    \begin{alertblock}{Constat}
        L'autonomie est divisée par 3 entre la surface et 20m !
        Plus on va profond, plus l'air se consomme vite.
    \end{alertblock}
\end{frame}

\begin{frame}{Astuce Pratique : La consommation au Palier (3m)}
    Comment estimer simplement sa consommation en bars pendant le palier ?
    
    \begin{block}{L'Astuce "Facile" (Règle du pouce)}
        Pour simplifier les calculs mentaux sous l'eau, retenez :
        \centering\Large\textbf{\textcolor{ffessmred}{2 bars par minute}}
    \end{block}

    \textbf{La Logique (Approximative) :}
    \begin{itemize}
        \item Consommation moyenne surface $\approx$ 20 L/min.
        \item À 3m (1,3 bar) $\rightarrow$ $20 \times 1,3 = 26$ L/min.
        \item \textbf{Bloc 15L} : $26 \div 15 \approx \mathbf{1,7}$ bar/min $\rightarrow$ Arrondi à \textbf{2 bars}.
        \item \textbf{Bloc 12L} : $26 \div 12 \approx \mathbf{2,1}$ bars/min $\rightarrow$ Arrondi à \textbf{2 bars}.
    \end{itemize}

    \begin{exampleblock}{Exemple concret}
        Pour un palier de \textbf{3 minutes} :
        $ 3 \text{ min} \times 2 \text{ bars/min} = \mathbf{6 \text{ bars consommés}} $
    \end{exampleblock}
\end{frame}

\begin{frame}{Planification et Gestion de l'Air (1/2)}

    % 1. L'Objectif de sécurité
    \begin{alertblock}{Objectif Prioritaire}
        En autonomie, la planification est indispensable pour \textbf{éviter la panne d'air}.
    \end{alertblock}

    \vspace{0.4cm}

    % 2. Les paramètres du calcul
    \textbf{Les 3 paramètres à prendre en compte :}
    \begin{enumerate}
        \item \textbf{Le Profil de la plongée} : La profondeur influe directement sur la consommation.
        \item \textbf{Le Volume de la bouteille} : La quantité d'air disponible au départ.
        \item \textbf{La Consommation du plongeur} : Un paramètre personnel et variable.
    \end{enumerate}

    \vspace{0.3cm}
\end{frame}

\begin{frame}{Planification et Gestion de l'Air (2/2)}
    % 3. La variabilité de la consommation (Note importante)
    \begin{block}{Attention à la consommation variable}
        Ce paramètre est difficile à évaluer car il évolue au cours de la plongée selon :
        \begin{itemize}
            \item \textbf{Le Froid}
            \item \textbf{Les Efforts} (courant, palmage, stabilisation)
            \item Le Stress
        \end{itemize}
    \end{block}

\end{frame}

\begin{frame}{Planification : La Règle des Tiers}
    Pour gérer sa sécurité et celle de la palanquée, on utilise souvent la règle des tiers pour gérer le gaz.

    \begin{columns}
        \column{0.5\textwidth}
        \begin{itemize}
            \item \textbf{1/3} : Pour l'aller / le fond.
            \item \textbf{1/3} : Pour le retour / la DTR.
            \item \textbf{1/3} : Sécurité (imprévus, assistance).
        \end{itemize}
       
        \column{0.5\textwidth}
        \begin{block}{Exemple (15L à 200b)}
            \begin{itemize}
                \item Air total : 3000 Litres.
                \item Air utilisable fond : 1000 Litres.
                \item À 40m (Conso 75L/min) :
                $$ 1000 / 75 \approx 13 \text{ min} $$
            \end{itemize}
        \end{block}
    \end{columns}

    \begin{alertblock}{Info Consommation}
        Si on n'a jamais mesuré sa consommation personnelle, on prend comme référence \textbf{15 L/min} (en surface).
    \end{alertblock}
\end{frame}
